% =========================================================
% Rational Approximation
% =========================================================

\section{Rational Approximation}
\label{sec:rational-approximation}

\subsection{Definitions and Theorems}

One of the most important structural facts about the rational numbers is that they can approximate
other numbers arbitrarily well. This is the point at which the density of $\mathbb{Q}$ begins to turn
into analysis.

\begin{theorem}[Rational approximation]
\label{thm:rational-approximation}
For every real number $x$ and every $\varepsilon > 0$, there exists $r \in \mathbb{Q}$ such that
\[
|x - r| < \varepsilon.
\]
\end{theorem}

\begin{remark*}[Proof]
\hyperref[prf:rational-approximation]{Go to proof.}
\end{remark*}

\begin{remark}
This theorem says that rational numbers are dense not only in the order-theoretic sense of lying
between real numbers, but also in the metric sense of lying arbitrarily close to any prescribed real
number.
\end{remark}

\begin{remark}[Constructive approximation methods]
There are several useful ways to approximate a real number by rational numbers.
\begin{enumerate}
    \item \textbf{Decimal truncation.} If $x = 1.4142135\dots$, then
    \[
    1.4, \quad 1.41, \quad 1.414, \quad 1.4142
    \]
    are rational approximations.
    \item \textbf{Mediants.} If
    \[
    \frac{a}{b} < x < \frac{c}{d},
    \]
    then the mediant
    \[
    \frac{a+c}{b+d}
    \]
    gives a new rational approximation between the two old ones.
    \item \textbf{Continued fractions.} Continued fractions often produce exceptionally good rational
    approximations and are closely related to the Stern--Brocot tree.
\end{enumerate}
\end{remark}

\begin{remark}[Approximation by sequences]
A real number can often be approximated naturally by a sequence of rationals. For example,
\[
1,\ 1.4,\ 1.41,\ 1.414,\ 1.4142,\dots
\]
is a sequence of rational approximations to $\sqrt{2}$. This point of view leads directly to Cauchy
sequences and to the construction of the real numbers.
\end{remark}

\subsection{Dirichlet's Approximation Theorem}

Density alone says that rationals come arbitrarily close to a real number, but it says nothing about
\emph{how efficient} those approximations can be made. Dirichlet's theorem gives a surprisingly strong
quantitative answer.

\begin{theorem}[Dirichlet's Approximation Theorem]
\label{thm:dirichlet-approx}
For every real number $x$ there exist infinitely many rational numbers $p/q$ with $p \in \mathbb{Z}$ and
$q \in \mathbb{N}$ such that
\[
\left|x - \frac{p}{q}\right| < \frac{1}{q^2}.
\]
\end{theorem}

\begin{remark*}[Proof]
\hyperref[prf:dirichlet-approx]{Go to proof.}
\end{remark*}

\begin{remark}[Interpretation]
If the denominator $q$ is moderately large, the error bound $1/q^2$ can already be very small. In
this sense the rational numbers approximate real numbers far better than density by itself might
suggest.
\end{remark}

\begin{remark}[Idea of proof]
Consider the fractional parts of
\[
x, 2x, 3x, \dots, (N+1)x.
\]
These $N+1$ numbers lie in the interval $[0,1)$, which is divided into $N$ equal subintervals.
By the pigeonhole principle, two of the fractional parts must land in the same subinterval. Their
difference produces integers $p$ and $q$ for which
\[
\left|x - \frac{p}{q}\right| < \frac{1}{qN}.
\]
A suitable choice of $N$ then yields the stated estimate.
\end{remark}

\begin{remark}[Conceptual significance]
The study of how well real numbers can be approximated by rationals is called
\emph{Diophantine approximation}. Dirichlet's theorem is one of its basic foundational results and
forms a natural bridge from elementary rational arithmetic to deeper number theory.
\end{remark}

% =========================================================
% Exposition: The Geometry of Rational Approximation
% =========================================================

\newpage
\section*{Reading the Diagram: A Walking Tour}

The figure illustrates three layers, each telling the same story from a different vantage point.
Take a moment to orient yourself before reading further.

\subsection*{Why rational approximation matters}

You have just proved that the rationals are dense in $\mathbb{R}$: between any two real numbers lies a rational one. That result tells you approximation is \emph{possible}, but it says nothing about how \emph{efficient} those approximations can be made.

The sharper question: given an irrational number $x$, how close can a fraction $p/q$ get to $x$, and how does the quality of the approximation depend on the size of the denominator $q$? 

\subsection*{Layer 1: The Farey gap identity}

The number line at the top shows two fractions in \emph{lowest terms}, $a/b$ and $c/d$, that are \emph{Farey neighbors}: $|bc - ad| = 1$. The blue tick marks the \emph{mediant} $(a+c)/(b+d)$, which lies strictly between them.

The key identity is exact:
\[
\frac{c}{d} - \frac{a}{b} = \frac{1}{bd}.
\]
The mediant splits the interval into two pieces of \emph{unequal} size:
\[
\frac{a+c}{b+d} - \frac{a}{b} = \frac{1}{b(b+d)}, \qquad \frac{c}{d} - \frac{a+c}{b+d} = \frac{1}{d(b+d)}.
\]
The larger denominator produces the smaller subgap. This asymmetry is the engine of the whole theory: inserting mediants refines intervals unevenly, always shrinking more on the side with the larger existing denominator.

\subsection*{Layer 2: The Stern--Brocot tree}

Starting from $0/1$ and $1/0$, we generate the \emph{Stern--Brocot tree}.
\begin{itemize}
    \item \textbf{Level 0:} Seeds $0/1, 1/1, 1/0$.
    \item \textbf{Level 1:} Mediants yield $1/2$ and $2/1$.
    \item \textbf{Level 2:} Yields $1/3, 2/3, 3/2, 3/1$.
\end{itemize}
Every positive rational appears in the tree exactly once in lowest terms. The tree encodes that simple fractions (small denominators) are coarse approximations, while complex fractions (large denominators) are needed for high precision.

\subsection*{Layer 3: Dirichlet approximation}

The bottom number line shows an irrational $x$ (red tick). Around it float four rational approximants $p_i/q_i$. The dashed circles represent the Dirichlet bound:
\[
\left| x - \frac{p}{q} \right| < \frac{1}{q^2}.
\]
As $q$ grows, the window $1/q^2$ tightens around $x$. Since $x$ must land in a Farey interval, the nearer neighbor satisfies this $1/q^2$ bound.

\subsection*{The Hidden Unity}

The invariant connecting all three panels is the determinant $|bc - ad| = 1$.
\begin{itemize}
    \item \textbf{Layer 1:} It gives the exact gap $1/(bd)$.
    \item \textbf{Layer 2:} It is preserved by mediant insertion.
    \item \textbf{Layer 3:} It implies the Dirichlet bound.
\end{itemize}
This identity is the thread connecting rational arithmetic to the modular group $\mathrm{SL}(2,\mathbb{Z})$.